\scanpage{35}
\mspage{17}
inverse de $x$ dans le ``rev. universel abélien'' $\widetilde G$ de $G$, construit comme $\varprojlim$ des revêtements $G(n)=G$ de $G$, donnés par $n\operatorname{id}$. L'énoncé dit que si ce torseur est trivial -- i.e. si $\forall n$, $x$ se remonte -- alors $x=0$. C'est vrai en car. $p>0$ également -- mais dans ce cas les ${}_nG$ ne sont pas étales.

\margin{\deleted{\unsure{Le théorème} \illegible{2} \unsure{Attention} \illegible{2} \unsure{de cette question} l'\unsure{affirmation} \illegible{3} \unsure{fausse}...}}

Si cependant on remplace les ${}_{p^\nu}G$ par leurs quotients par leurs parties radicielles -- en travaillant donc dans la cat. des ``quasi-variétés'' où on \deleted{identifie} \inserted{inverse} les \deleted{\unsure{immersions|isomorphismes}} \inserted{morphismes} finis radiciels surjectifs, \deleted{\unsure{pour peu qu'il}} \inserted{pour travailler} avec la partie proétale de $H_1(\overline G)$, \deleted{\unsure{quitte}} l'existence de relèvements radiciels $/K$ implique déjà $x=0$. \inserted{Bien sûr $x$ est de $p$-torsion,} On est ramené au cas d'une VA, avec $px=0$.... Mais \deleted{la} validité sous cette forme du th. 3 \unsure{ou|en} \inserted{de son cor. en} car. $p>0$ nous fait une belle jambe -- car pour \unsure{être} à \unsure{l'aise|l'air}, il faudrait pouvoir le prouver \deleted{pour} \inserted{dans} les \unsure{contextes} des groupes fondamentaux ``premiers à $p$''. Or dans ce con\inserted{texte}, déjà dans le cas d'une courbe elliptique $Y$ (\deleted{\unsure{courbe}} \inserted{rendue} anabélienne au besoin en enlevant un pt) l'injectivité de (**) est \uline{fausse}. Il

\margin{O.K. en se remplaçant d'abord $K$ par un \unsure{alg.|anneau} de type fini sur $\mathbb F_p$, puis en \unsure{se ramenant} \unsure{au cas} d'un corps fini \uline{parfait}.}
\margin{34}
